%%%%%%%%%
\section{Graded unitarity constraints on $V_k(\mathfrak{sl}_N)$} \label{sec:typeAgrunit}
%%%%%%%%%
In this section, we investigate constraints imposed by graded unitarity on $V_k(\mathfrak{sl}_N)$ affine Kac--Moody algebras. These constraints are implemented at the level of the vacuum determinant in a fixed-weight sector, and hence are weaker than the full set of graded unitarity constraints on the vertex algebra. 

Such an analysis has already been carried out for $\gf = \slf_2, \slf_3$ and $\slf_4$ \cite{ArabiArdehali:2025fad}, where it was shown that graded unitarity constraints are very restrictive and rule out all levels except,
\begin{equation}
\begin{split}
    &k = k_{h^\vee,q} = -h^\vee + \frac{h^\vee}{q}: (h^\vee,q) = 1, \qquad \mathrm{for} \ \ \gf = \mathfrak{sl}_2, \mathfrak{sl}_3~, \\
    &k = k_{4,q} = -4 + \frac{4}{q} \ \ \ \ \ :4 \nmid q, \qquad \qquad \ \ \ \mathrm{for} \ \ \gf = \mathfrak{sl}_4~.
\end{split}
\end{equation}
Note that for $\gf = \mathfrak{sl}_4$, there is no coprimality condition and it was shown \emph{loc. cit.} that, in addition to the boundary admissible levels, fixed root determinant-based constraints do not rule out the non-admissible $k = k_{2,q}$ levels with $q$ odd. We show that this pattern generalises to $\gf = \mathfrak{sl}_N$ for all $N$. 

In what follows, we restrict our attention exclusively to $\gf = \slf_N$ with $N \geqslant 5$. The positive roots and the Weyl vector for $\slf_N$ are given by,
\begin{equation}
\begin{split}
    &\Delta_+ = \{\alpha_{r,s} := e_r - e_s: 1 \leqslant r < s \leqslant N \}, \\
    &\rho = \frac{1}{2} \sum_{\alpha \in \Delta_+} \alpha = \left(\frac{N-1}{2}~, \frac{N-3}{2}, \ldots, -\frac{(N-1)}{2} \right)~.
\end{split}
\end{equation}
For our purposes, it is relevant to simplify the regularity criterion \eqref{eq:ADEregularity} into something more amenable to calculations in $\gf = \slf_N$. The weights $\lambda$ appearing in the argument of $E(\lambda)$ in the exponent $M_{p/q}$ are always positive integral multiples of positive roots. Consequently, it can be shown by direct calculation that for $\gf = \slf_N$ and $\lambda = t \alpha_{r,s}$ (where $t \in \mathbb{Z}_{\geqslant 1}$ and $\alpha_{r,s} \in \Delta_+$), the regularity criterion simplifies to,\footnote{This result is in agreement with the simplification obtained in \cite{jiang2025singular}.}
\begin{equation} \label{eq:typeAregularity}
    E(t \alpha_{r,s}) \neq 0 \Longleftrightarrow t \geqslant r \ \ \mathrm{and} \ \ t \geqslant N - s+1~.
\end{equation}
The approach to implementing graded unitarity constraints is analogous to \cite{ArabiArdehali:2025fad}. For a given conformal level $L$, we first determine the smallest values of $A_{(t,\alpha)} := t(t+(\rho|\alpha))$ such that $E(t\alpha) \neq 0$ in \eqref{eq:splitM}. Since $\frac{A_{(t,\alpha)}}{L} \leqslant \frac{p}{q}$, the smallest values of $A_{(t,\alpha)}$ determine the most negative and the second most negative zero of the vacuum determinant in a fixed weight $\nu = \alpha$ sector, where $\alpha$ is a root of $\slf_N$. The multiplicities of these zeros are then found using \eqref{eq:expinchar}.

To this end, note that the simplified $\slf_N$ regularity criterion \eqref{eq:typeAregularity} implies that $(\rho|\alpha) \geqslant N - 2t + 1$. Consequently, one has $A_{(t,\alpha)} \geqslant t(N-t+1)$, leading to,
\begin{equation}
\begin{split}
    A_{(t,\alpha)} \begin{cases}
        = N \qquad \qquad \ \ \ \mathrm{when} \ \ t = 1 \ \ \mathrm{and} \ \ \alpha = \theta \\
        \geqslant (2N-2) \qquad \mathrm{otherwise}
    \end{cases}
\end{split}
\end{equation}
Thus, when $N \nmid L$, the $k = k_{N,L}$ zero in the fixed-root vacuum determinant at conformal level $L$ is a consequence of the primitive null state in the adjoint representation. The multiplicity of this zero in the vacuum determinant $\mathrm{det}~S_{L \delta + \alpha}$ can be computed from,
\begin{equation}
    \frac{M_{N/L,L}}{E(0)} = \frac{E(\theta)}{E(0)} = \mathrm{ch~Ad}~,
\end{equation}
to be equal to one. It is worth stressing that this result does \emph{not} depend on any coprimality condition between $N$ and $L$. To summarise, the non-critical zeros of the fixed root, conformal level $L$ vacuum determinant det $S_{L \delta + \alpha}$ are given by,
\begin{equation} \label{eq: typeAprelimzero}
    k = \begin{cases}
        k_{N,L} = -N + \frac{N}{L}: N \nmid L \\ 
        k_{N,q} = -N + \frac{N}{q}: q < L \ \ \mathrm{and} \ \ N \nmid q \\
        k_{A,B} = -N + \frac{A}{B}: A \geqslant 2N-2 \ \ \mathrm{and} \ \ B \leqslant L
    \end{cases}
\end{equation}

\subsection{Structure of zeros of the vacuum determinant}
Before discussing the precise structure of zeros of the fixed-root vacuum determinant, it is worthwhile to perform a bespoke analysis at $L = 1$ and $L = 2$. Note that the fixed-root vacuum determinant for $\gf = \slf_N$ at these levels is given by,
\begin{equation} \label{eq:lowleveltypeAdet}
\begin{split}
    &\mathrm{det}~S_{\delta + \alpha} \propto (k-k_{N,1}) \\
    &\qquad \qquad \propto \textcolor{blue}{(k - 0)}~, \\
    &\mathrm{det}~S_{2\delta + \alpha} \propto (k - k_{N,2})(k-k_{N-1,1})^{N-3} (k - k_{N,1})^{m_{N/1}(2\delta + \alpha)} (k - k_{N+1,1})^{m_{(N+1)/1}(2\delta + \alpha)}\\ 
    &\qquad \qquad  \propto \textcolor{blue}{\left(k + \frac{N}{2} \right)(k + 1)^{N-3}} (k - 0)^{m_{N/1}(2\delta + \alpha)} (k - 1)^{m_{(N+1)/1}(2\delta + \alpha)}~.
\end{split}
\end{equation}
Hence, from graded unitarity constraints at level one, and from the fact that the Segal--Sugawara stress tensor has $R$-charge $R(T) = 1$, one obtains that $-N < k \leqslant 0$. Furthermore, graded unitarity constraints at level two additionally rule out $k \in \left(-\frac{N}{2},-1\right)$. The $k = -N/2$ zero is a consequence of a primitive null state in the adjoint representation, while the zero at $k = -1$ comes from a state in the $[0,1,0,\ldots,1,0]$ representation. Passing to the adjoint representation basis instead of a fixed weight basis, the norm sign incompatibility with graded unitarity, of the state in the adjoint representation persists in the whole regime $k \in \left(-\frac{N}{2}, 0\right)$. Consequently, at level two, constraints from graded unitarity restrict the level $k$ to be in the regime $-N < k \leqslant -\frac{N}{2}$ or $k = 0$, where the latter corresponds to the trivial vertex algebra.

To implement graded unitarity constraints for $L \geqslant 3$, we will use \eqref{eq: typeAprelimzero} to find the most negative and the second most negative zeros in the vacuum determinant in the fixed-root sector. Consider a conformal level $L$ such that $N \nmid L$. From \eqref{eq: typeAprelimzero}, the most negative zero of the vacuum determinant is $k = k_{N,L}$. To find the second most negative zero, we divide the analysis into two cases,
\begin{itemize}
    \item Case 1: $N \nmid (L-1)$ \\ 
    \\
    Since $N \nmid (L-1)$, $k = k_{N,L-1}$ first appears as a zero of the vacuum determinant at conformal level $L-1$ with multiplicity one, and hence, is also a zero of the vacuum determinant at level $L$. For $N \geqslant 5$ and  $L \geqslant 3$, one obtains,
    \begin{equation} \label{eq:typeA2ndsmallest}
    \begin{split}
        &(2N-2)(L-1) - NL = (N-2)L - (2N-2) > 0 \\
        &\Longrightarrow k_{2N-2,L} = -N + \frac{2N-2}{L} > -N + \frac{N}{L-1} = k_{N,L-1}~.
    \end{split}
    \end{equation}
    This shows that the zeros of the third kind in \eqref{eq: typeAprelimzero} are always greater than $k = k_{N,L-1}$ and hence for conformal levels $L$ such that $N \nmid (L-1)$, the second most negative zero is $k = k_{N,L-1}$. 
    \item Case 2: $N | (L-1)$ \\ 
    \\
    In this case, $N \nmid (L-2)$. Thus, $k = k_{N,L-2}$ first appears as a zero of the vacuum determinant at conformal level $L-2$ with multiplicity one, and hence, is also a zero of the vacuum determinant at level $L$. For $N \geqslant 5$ and $N | (L-1)$, one obtains,
    \begin{equation}
    \begin{split}
        &(2N-2)(L-2) - NL = (N-2)L - (4N-4) > 0 \\
        &\Longrightarrow k_{2N-2,L} = -N + \frac{2N-2}{L} > -N + \frac{N}{L-2} = k_{N,L-2}~.
    \end{split}
    \end{equation}
    This shows that the second most negative zero of the vacuum determinant is $k = k_{N, L-2}$. 
\end{itemize}
Finally, for $L$ that is a multiple of $N$, \emph{i.e.}, $N | L$, the most negative zero is given by $k = k_{N, L-1}$ by an argument analogous to the one leading to \eqref{eq:typeA2ndsmallest}. 

To summarise, we obtain the following structure of zeros for $N \geqslant 5$ and $L \geqslant 3$, 
\begin{itemize}
    \item At level $L \neq (0,1)~\mathrm{mod}~N$, the most negative zero of the fixed-root vacuum determinant is given by $k = k_{N,L}$ and this zero always occurs with multiplicity one. Moreover, the second most negative zero is given by $k = k_{N,L-1}$. 
    \item At level $L = 0~\mathrm{mod}~N$, the most negative zero of the fixed-root vacuum determinant is given by $k = k_{N,L-1}$. 
    \item At level $L = 1~\mathrm{mod}~N$, the most negative zero of the fixed-root vacuum determinant is given by $k = k_{N,L}$ and this zero always occurs with multiplicity one. Moreover, the second most negative zero is given by $k = k_{N,L-2}$. 
\end{itemize}

\subsection{Constraints on the level}
The conclusion of the structure of zeros of the previous subsection is that the fixed-root vacuum determinant at level $L \geqslant 3$ takes the following form, \\
\\
\underline{\textbf{Case 1: } $L \neq 0~\mathrm{mod}~N$}
\begin{equation} \label{eq:typeAvacdet1}
\begin{split}
    &\mathrm{det}~S_{L\delta+\alpha} \propto \textcolor{blue}{(k-k_{N,L})^1}(k-k_{N,L-1})^{\#_1}( \ldots \ldots), \qquad L \neq 1~\mathrm{mod}~N~, \\
    &\mathrm{det}~S_{L\delta+\alpha} \propto \textcolor{blue}{(k-k_{N,L})^1}(k-k_{N,L-2})^{\#_2}( \ldots \ldots), \qquad L = 1~\mathrm{mod}~N~,
\end{split}
\end{equation}
where $( \ldots \ldots)$ denote other zeros of the vacuum determinant that are strictly greater (less negative) than $k = k_{N,L-1}$ (or $k_{N,L-2}$). $\#_1$ and $\#_2$ are some positive integers, and the proportionality factor is a $k$-independent rational number. 
\\
\\
\underline{\textbf{Case 2: } $L = 0~\mathrm{mod}~N$}
\begin{equation} \label{eq:typeAvacdet2}
    \mathrm{det}~S_{L\delta+\alpha} \propto \textcolor{blue}{(k-k_{N,L-1})^{\#_3}}( \ldots \ldots)~,
\end{equation}
where again $( \ldots \ldots)$ indicate zeros strictly greater than $k = k_{N,L-1}$ and $\#_3$ is some positive integer.\footnote{Throughout the paper, we adopt the convention that any omitted zeros in the relevant vacuum determinant are strictly greater than those displayed, $\#$ denotes a positive integer, and the proportionality factors are positive.} 

The all-level determinant-based constraints from graded unitarity now follow almost immediately. In particular, let us first consider $L \neq 0~\mathrm{mod}~N$. From \eqref{eq:typeAvacdet1}, we see that the vacuum determinant flips sign at $k = k_{N,L}$, with the next zero being at $k_{N,L-1}$ or $k_{N,L-2}$. Hence, at level $L \neq 0~\mathrm{mod}~N$, the graded unitarity sign rule \eqref{eq:grunitsignrule} eliminates the regime $k \in (k_{N,L}, k_{N,L-1})$ or $k \in (k_{N,L}, k_{N,L-2})$ (depending on whether $L \neq 1~\mathrm{mod}~N$ or $L = 1~\mathrm{mod}~N$), respectively. 

Turning to the $L = 0~\mathrm{mod}~N$ case, note that the most negative zero is $k = k_{N,L-1}$, and all the $k_{N,q}$ zeros (with $1 < q \leqslant L-1$ and $N \nmid q$) of the level $L-1$ vacuum determinant are also zeros of the level $L$ vacuum determinant. Since $L-1 \neq 1~\mathrm{mod}~N$, graded unitarity constraints already rule out the region $(k_{N,L-1}, k_{N,L-2})$ at conformal level $L-1$. Thus, we do not get any new constraints from graded unitarity at level $L = 0~\mathrm{mod}~N$. 

Therefore, we find that with the assumption of the standard $\mfR$-filtration on the $V_k(\mathfrak{sl}_N)$ affine Kac--Moody algebras, the all-level (fixed root) determinant-based graded unitarity constraints rule out all values of the level $k$ other than,
\begin{equation} \label{eq:typeAalllevel}
    k = \begin{cases}
        \textcolor{knownSCFTtext}{k_{N,q} = -N + \frac{N}{q}: (N,q) = 1} \\
        k_{N,q} = -N + \frac{N}{q}: (N,q) > 1 \ \ \mathrm{and} \ \ N \nmid q
    \end{cases}
\end{equation}
These constraints are a natural generalisation of the $\gf = \mathfrak{sl}_2, \mathfrak{sl}_3$, and $\mathfrak{sl}_4$ exclusions of \cite{ArabiArdehali:2025fad}, and demonstrate that the determinant-based constraints are independent of whether $N$ and $q$ are coprime. The first set of levels above is the set of boundary admissible levels. The $V_k(\slf_N)$ affine Kac--Moody algebras corresponding to these levels are precisely the ones associated with the ``class I" $D_q^N(\mathfrak{sl}_N)$ Argyres--Douglas SCFTs under the SCFT/VOA correspondence \cite{Xie:2016evu, Song:2017oew}. 

On the other hand, the second set of levels are all non-admissible. These appear as the levels of the $\mathfrak{sl}_N$ affine Kac--Moody subalgebras of the associated VOAs of $D_q^N(\mathfrak{sl}_N)$ Argyres--Douglas SCFTs having extra mass parameters/flavour symmetries coming from the irregular puncture; hence the VOAs associated with these theories are strictly \emph{larger} than $V_k(\mathfrak{sl}_N)$ \cite{Xie:2016evu, Song:2017oew, Beem:2023ofp}. While allowed by all-level (fixed root) determinant-based graded unitarity constraints (under the assumption of the standard $\mathfrak{R}$-filtration), we leave open the possibility that an analogous analysis in a different weight sector or a finer state-based graded unitarity analysis might rule them out. 

As an implication of \eqref{eq:typeAalllevel}, note that when $N$ is prime, determinant-based graded unitarity constraints are sufficient to rule out \textit{all} levels other than the boundary admissible $k = k_{N,q}$, precisely the ones associated with the ``class I" $D_q^N(\mathfrak{sl}_N)$ Argyres--Douglas SCFTs. 